\chapter{\MakeUppercase{Introduction}}
\justifying
\section{Background}
\acp{dtn} are networks in which a continuous path between source and destination cannot be assumed. Nodes store messages, carry them while moving, and forward them opportunistically when contacts occur. This store-carry-forward model is useful in disaster-response scenarios because damaged infrastructure, sparse deployment, intermittent wireless links, and node mobility make conventional Internet-style routing impractical \citep{cerf2007dtn}. In such environments, communication support must continue even when connectivity is fragmented.

The project studies routing for a disaster-response \ac{dtn} consisting of civilians, responders, shelters, and drones. Each role has different mobility and operational meaning. Civilians move slowly and may carry local reports. Responders move faster and are expected to participate in rescue operations. Shelters are static destinations or coordination points. Drones are fast mobile relays that can bridge separated areas. A routing protocol that treats all of these nodes identically ignores information that is already available in the scenario.

The protocol proposed in this thesis is called Spatio-Semantic routing. The term spatial refers to location, zone visitation, and physical proximity to the destination. The term semantic refers to the operational role of the node and the priority of the message. The protocol combines both dimensions with encounter history and transitive delivery prediction. Its aim is not simply to maximize all message delivery, but to preserve the delivery of critical messages when traffic load increases and buffers become congested.

\section{Need for Delay Tolerant Routing}
Traditional \ac{manet} routing protocols generally depend on the existence or discovery of a contemporaneous path. This assumption is weak in sparse, mobile, and damaged-network conditions. In contrast, \ac{dtn} routing accepts that contacts may be rare and that the best relay at the present time may only be useful because of where it is likely to move later. Routing therefore becomes a decision problem: when two nodes meet, each protocol must decide which messages should be copied, how many copies should be transferred, and whether the receiver is more useful than the current carrier.

Three broad routing behaviors are relevant to this project. Epidemic routing floods messages to every encountered node that does not already hold the message \citep{vahdat2000epidemic}. Spray and Wait limits the number of copies and then waits for a direct encounter with the destination \citep{spyropoulos2005spray}. Prediction-based protocols such as \ac{prophet} use encounter history and transitivity to forward to more promising carriers \citep{lindgren2012prophet}. The proposed protocol draws from these ideas but adds spatial zones, semantic node roles, critical-message buffer reservation, and hybrid copy allocation.

\section{Motivation}
The practical motivation is emergency communication. In a disaster area, messages are not equally important. A medical alert, shelter request, or responder coordination update should not be treated like routine background information. A routing protocol that achieves high average delivery but loses urgent messages under congestion is unsuitable for this use case.

The technical motivation is that existing baselines expose opposite weaknesses. Epidemic routing can deliver well in light traffic, but its uncontrolled replication creates high overhead and buffer pressure. Spray and Wait controls overhead, but its conservative forwarding can delay or miss urgent messages when the limited copies are held by weak carriers. The proposed protocol attempts to occupy the useful middle ground: it forwards aggressively when a message is critical or when enough copies exist, but it becomes selective as copies become scarce.

\section{Disaster-Response Communication Characteristics}
Disaster communication differs from ordinary mobile networking in several important ways. First, the physical environment is unstable. Roads may be blocked, electricity may be unavailable, and parts of the affected area may become temporarily isolated. A message may therefore depend on a sequence of opportunistic contacts across civilians, responders, or drones instead of a stable wireless infrastructure. Routing decisions made under these conditions must tolerate uncertainty rather than assume that a short path currently exists.

Second, the participating nodes do not behave uniformly. Civilians often move unpredictably and may remain within a local neighborhood while seeking safety or assistance. Responders typically move with clearer operational intent and may revisit command areas or incident zones repeatedly. Shelters remain fixed and can act as communication anchors. Drones provide rapid movement across the map and can bridge disconnected regions that are difficult for ground carriers to connect. Because these roles imply different future forwarding value, role-aware routing is a natural extension rather than an artificial feature.

Third, the traffic itself is heterogeneous. Some messages are delay sensitive because they represent injury reports, requests for evacuation, or updates that affect ongoing rescue actions. Other messages are still useful but less urgent. A protocol that treats every packet identically ignores the actual purpose of communication in disaster operations. This thesis therefore treats message priority as part of routing logic instead of leaving it outside the forwarding process.

Resource scarcity is another defining characteristic. Buffer capacity, wireless contact duration, and useful transmission opportunities are all limited. In \acp{dtn}, an incorrect forwarding decision has a long-lasting effect because a wasted copy may remain unavailable until the next rare contact. For this reason, disaster routing should not only maximize the number of deliveries. It should also improve the quality of those deliveries by protecting critical traffic and choosing carriers with stronger future usefulness.

\section{Scope of the Project}
The scope of the project includes the design, implementation, and simulation-based evaluation of a spatio-semantic routing protocol for a disaster-response \ac{dtn}. The implementation is written in Python and compares three protocols:

\begin{itemize}
	\item Epidemic routing as the flooding baseline.
	\item Spray and Wait routing as the copy-limited baseline.
	\item Spatio-Semantic routing as the proposed protocol.
\end{itemize}

The simulator models a square disaster area, wireless contact range, finite node buffers, message generation under three traffic levels, critical and non-critical messages, random waypoint mobility, static shelters, and fast drones. The evaluation uses delivery ratio, critical delivery ratio, average delay, average critical delay, overhead ratio, and buffer drops.

The project does not implement a real wireless deployment, cryptographic security, energy consumption, interference modeling, or real map-based mobility. These are outside the current thesis scope. The results should therefore be interpreted as simulation evidence for the protocol design rather than as direct field deployment measurements.

\section{Assumptions and Thesis Organization}
The simulator adopts a deliberately abstract representation of the disaster area. Nodes move in continuous two-dimensional space, contacts are determined by geometric range, and message exchanges occur without detailed radio-layer contention. These assumptions are useful because they allow the thesis to isolate the routing question: given limited contacts and limited storage, which forwarding strategy is most effective? The abstraction therefore reduces implementation noise while preserving the central delivery trade-offs of a sparse mobile network.

This choice is also important for academic clarity. A disaster network can be modeled at many levels, from radio propagation to social behavior to vehicle-road interaction. Including all of them in one project would produce a broad but less focused study. The present thesis instead concentrates on forwarding intelligence under disruption. By doing so, it becomes possible to examine how routing changes when the protocol is aware of who is carrying the message, where that carrier tends to move, and how urgent the message is.

The thesis is organized to move from motivation to evidence. Chapter 2 reviews relevant routing literature, identifies the disaster-response research gap, and defines the project goals. Chapter 3 formalizes the technical specification, simulation requirements, and evaluation parameters. Chapter 4 describes the overall system design and the internal logic of the proposed Spatio-Semantic router. Chapter 5 explains the testing methodology, traffic scenarios, and success criteria used for comparison. Chapter 6 presents the implementation details of the simulator, routing modules, and visualization tools. Chapter 7 analyzes the measured results across traffic levels and interprets the trade-offs among delivery, delay, overhead, and drops. Finally, Chapter 8 summarizes the contributions of the work and outlines directions for further enhancement.

Taken together, these chapters support a single argument. Opportunistic forwarding in a disaster \ac{dtn} should not be blind, purely probabilistic, or indifferent to message importance. It should use available contextual signals to improve the treatment of urgent communication while still respecting the limited resources of a disrupted network. The remainder of the thesis develops and evaluates that argument in a systematic way.
